\chapter{What bag generate writes}
\label{ch:outputs}

\cmd{bag generate} reads a spec, checks it, and writes every generated file into one
folder, the \term{output directory}. You name it with \opt{-o}. Without \opt{-o}, BAG
uses Project > \setting{Output directory}, relative to the folder of the spec. This
chapter describes every file and folder that \cmd{bag generate} writes there, and what
each one is for.

The examples in this chapter use a small design, \code{lite_ss}: one AXI4-Lite master
\code{cpu}, an AXI4-Lite shared bus \code{bus}, and two AXI4-Lite slaves, \code{sram} and
\code{uart}, in one clock domain. Its IR, \file{lite_ss.ir.yaml}, is in \file{files/} of
the examples. \cmd{bag generate} accepts an IR in place of a workbook (chapter~\ref{ch:ir}).

\section{The output tree}
\label{sec:tree}

\lstinputlisting[style=console]{examples/console/out-tree.txt}

The summary line names the \term{top}: the module that the lint, compile, simulation and
formal steps use as their top module. In the flat port style it is the top wrapper.

\begin{description}
\item[\file{rtl/}] the SystemVerilog of the design (section~\ref{sec:rtl}).
\item[\file{lite_ss.f}] the file list, \file{<project>.f} (section~\ref{sec:filelist}).
\item[\file{sw/}] the C header (section~\ref{sec:cheader}).
\item[\file{doc/}] the documents (section~\ref{sec:docs}).
\item[\file{tests/}] the generated cocotb tests (section~\ref{sec:tests}).
\item[\file{formal/}] the SymbiYosys jobs and their properties (section~\ref{sec:jobs}).
\item[\file{bag_outputs.json}] the manifest: what this run wrote
  (section~\ref{sec:manifest}).
\end{description}

Which of these parts BAG writes depends on settings of the Project sheet
(section~\ref{sec:switches}). The commands \cmd{bag lint}, \cmd{bag sim} and
\cmd{bag formal} write more folders into the output directory when they run
(chapter~\ref{ch:flows}). \cmd{bag generate} does not write them.

\section{The RTL}
\label{sec:rtl}

Everything under \file{rtl/} is SystemVerilog (IEEE 1800-2017).

\begin{description}
\item[\file{rtl/pkg/<package>.sv}] The \term{parameter package}. Its name is Wrapper >
  \setting{Parameter package}. It holds, for each fabric, the widths, port counts and
  address decode tables that the fabric's template reads, and the memory map as
  constants: \code{<Slave><Region>Base} and \code{<Slave><Region>Size}, in system
  addresses. Your SoC can import the package to use the same addresses as the fabric.
\item[\file{rtl/<top>.sv}] The \term{top wrapper}. Its name is Wrapper >
  \setting{Top module name}. It is the module that your SoC instantiates. It has its
  clock and reset ports (Wrapper > \setting{Clock/reset ports}), and one set of bus
  signals per master and slave port, named by Wrapper > \setting{Signal naming pattern}.
  Inside, it conditions the resets that need it (an in-house \code{bag_rst_sync} module, see
  \file{doc/reset.md}) and instantiates the fabrics, bridges and adapters.
\item[\file{rtl/<sub-wrapper>.sv}] One module for each sub-wrapper of the Hierarchy
  sheet. With a \setting{Module prefix} the module and the file are \code{<prefix>_<name>}
  (section~\ref{sec:prefix}). A sub-wrapper groups some instances inside the top. Its ports are always flat,
  one port per signal, whatever the port style.
\item[\file{rtl/inhouse/}] Copies of the in-house RTL that the design uses: the modules
  written for BAG, such as the AHB-Lite and APB fabrics, the bridges between protocol
  families, the address decoder and the reset synchronizer. BAG copies only the files
  that the design needs.
\item[\file{rtl/vendor/pulp/}] Copies of the files of the PULP \code{axi} and
  \code{common_cells} libraries that the design uses, with the \file{LICENSE} file of
  each library. The include files are here too, under \file{include/}.
\item[\file{rtl/vendor/derived/}] Files that BAG derives from PULP files, such as
  \file{bag_axil_mux.sv}. Each one names its upstream file and commit, lists what was
  changed, and keeps the upstream licence notice.
\item[\file{rtl/ext/}] Only for a user-supplied fabric: a stub for you to replace
  (section~\ref{sec:external}).
\end{description}

With a blank \setting{Module prefix}, BAG copies the in-house and vendored files as they
are. It does not add its file header to them (section~\ref{sec:header}). The lint step treats every file under
\file{rtl/vendor/} as third-party code (section~\ref{sec:lint}).

\subsection{The module prefix}
\label{sec:prefix}

With a \setting{Module prefix}, for example \code{soc1}, BAG renames the in-house modules,
the modules it derives from PULP, and the sub-wrappers to \code{soc1_<name>}, in the
declaration, in every instantiation and in the \code{endmodule} label. The files of those
modules get the same prefix, and the file list follows. The top wrapper, the parameter
package, the interfaces and the flat harness keep their names, and so do instance names, so
hierarchical paths do not change. Comments and messages inside the files keep the old
names. A file that BAG derives from PULP also gets a line after its SPDX line that says
that its unit names carry the prefix. \file{doc/port_map.md} keeps the sub-wrapper names
of the Hierarchy sheet. The vendored PULP files are not renamed yet.

\lstinputlisting[style=console]{examples/console/out-prefix.txt}

\subsection{The two port styles}

Wrapper > \setting{Port style} sets the ports of the top wrapper. The option
\opt{--style} of \cmd{bag generate} overrides it for one run.

\begin{description}
\item[\code{flat}] One port per signal, for example \code{cpu_axil_awaddr}. This is the
  default. The top wrapper is also the top of lint, compile, simulation and formal.
\item[\code{sv-interface}] One SystemVerilog interface port per master and slave port.
  BAG writes two more files: \file{rtl/<project>_<proto>_if.sv}, one interface for each
  protocol at the top's ports (\code{axi}, \code{axil}, \code{ahb}, \code{apb3},
  \code{apb4}), and \file{rtl/<top>_flat.sv}, a \term{flat harness}. The harness has
  flat ports, instantiates the interfaces and the top wrapper, and connects them.
  Verilator cannot compile a top module with interface ports, so the harness is the top
  of lint, compile, simulation and formal.
\end{description}

\lstinputlisting[style=console]{examples/console/out-style.txt}

The summary line names \code{lite_ss_top_flat} as the top: the flat harness. Your SoC
still instantiates \code{lite_ss_top}.

\section{The file list}
\label{sec:filelist}

\file{<project>.f} lists the files of the design for your own tools. It starts with one
\code{+incdir+} line for each include directory that the vendored files need. Then it
names every file to compile, one per line, in compile order: the package first, then the
interfaces (\code{sv-interface} style only), the vendored and in-house files, the
sub-wrappers, the top wrapper and, last, the flat harness (\code{sv-interface} style
only). The paths are relative to the output directory, so run your tool from there.

\lstinputlisting[style=console]{examples/console/out-filelist.txt}

The \file{LICENSE} files and the include files are not in the list: they are not
compiled on their own.

\section{The manifest}
\label{sec:manifest}

\file{bag_outputs.json} records what \cmd{bag generate} wrote. \cmd{bag lint},
\cmd{bag sim} and \cmd{bag formal} read it to find the files, so they take only the
output directory as their argument. It is a JSON object with these keys:

\begin{description}
\item[\code{files}] every file written, relative to the output directory, with its
  SHA-256 hash.
\item[\code{top}] the top of lint, compile, simulation and formal.
\item[\code{filelist}] the name of the file list.
\item[\code{tests}] the generated cocotb module, or \code{null} if there is none.
\item[\code{formal}] the generated SymbiYosys jobs (\file{formal/*.sby}); empty if there
  are none.
\item[\code{lint_tools}] Project > \setting{Lint tools}.
\item[\code{simulator}] Project > \setting{Simulator}.
\end{description}

When you generate into the same output directory again, BAG reads the old manifest
first. It deletes each file that the old manifest lists and the new run does not write,
for example the flat harness after a change from \code{sv-interface} to \code{flat}.
BAG does not touch files that it did not write.

\section{The C header}
\label{sec:cheader}

\file{sw/<project>_memmap.h} gives software the memory map. For each region of the
\sheet{Memory_Map} sheet it defines \code{<SLAVE>_<REGION>_BASE} and \code{<SLAVE>_<REGION>_SIZE}
as \code{ULL} constants, in system addresses, grouped by fabric. A comment after each
base gives the access (\code{RW}, \code{RO} or \code{WO}) and marks a secure-only region
and a bridge window.

\lstinputlisting[style=console]{examples/console/out-c-header.txt}

\section{The documents}
\label{sec:docs}

\file{doc/} holds five documents for the people who integrate and review the design.

\lstinputlisting[style=console]{examples/console/out-docs.txt}

\begin{description}
\item[\file{doc/memmap.md} and \file{doc/memmap.html}] The memory map, per fabric: each
  region with its base, inclusive end, size, access and whether it is secure only; the
  bridge windows; each master's aliases; and what a master gets for an address that no
  region decodes. The HTML file shows the same rows as the Markdown file. It needs no
  other file, so you can open it in a browser or attach it to a review.
\item[\file{doc/port_map.md}] Per fabric: its protocol, topology, arbitration, template
  and instance name, and every master and slave port with its index, protocol, widths,
  clock domain, priority and the adapters on it. Use it to check how BAG resolved the
  ports.
\item[\file{doc/tieoffs.md}] The ports of BAG's components that are not simply wired to a
  bus signal: inputs tied to a constant, defaults, fitted widths and outputs that nothing
  reads, each with the reason. The report does not list every such port yet
  (chapter~\ref{ch:limits}).
\item[\file{doc/reset.md}] The \term{reset sequence} that the design needs: its clock
  domains and reset ports, every clock-domain crossing, the \term{reset groups} (the
  domains that must reset together because a crossing joins them), and the steps to
  follow when you assert and release the resets. It also says what goes wrong otherwise,
  and what the generated tests check of it. The file \file{docs/reset-sequence.md} in
  the BAG repository explains why.
\end{description}

\section{The generated tests}
\label{sec:tests}

\file{tests/} holds the cocotb tests for this design. \cmd{bag sim} runs them
(section~\ref{sec:sim}).

\begin{description}
\item[\file{tests/<project>_tests.py}] One cocotb test for each generated test that
  applies to this design. A test that could check nothing on this design is left out.
  The test bodies, the reference model and the bus models are part of BAG: the module
  imports them from \code{hdl_bag.tests_gen}.
\item[\file{tests/<project>.ir.yaml}] A copy of the IR. The tests read it to know what
  the design must do.
\item[\file{tests/<top>_ready_shim.sv}] The \term{READY shim}: a module with the ports of
  the top that instantiates the top and drives each AXI READY into the design from that
  channel's VALID. A design that waits for READY before it raises VALID then has a
  combinational loop, which does not settle in simulation. \cmd{bag sim}
  \opt{--ready-shim} uses it. BAG writes the shim only when the design has an AXI4 or
  AXI4-Lite port and the back-pressure test applies.
\end{description}

\lstinputlisting[style=console]{examples/console/out-tests-formal.txt}

\section{The formal jobs}
\label{sec:jobs}

\file{formal/} holds the SymbiYosys jobs and the properties they check. A \term{job} is
one \file{.sby} file. A job has one or more \term{tasks}; each task runs one check, in
one mode, to one depth (section~\ref{sec:formal}). Each job reads the design from
\file{rtl/} and a property file of its own. The property files attach to the design with
\code{bind}, so the design needs no formal wrapper.

\begin{description}
\item[\file{formal/<project>_reset.sv}, \file{.sby}] The quiet-in-reset job. While a
  domain's reset is held, the design asks nothing of the masters and slaves in that
  domain. Every design has this job.
\item[\file{formal/<project>_protocol.sv}, \file{.sby}] The port-protocol job. It checks
  each master and slave port of the top against the rules of its protocol. A design
  with no port to check has no protocol job.
\item[\file{formal/<project>_decode.sv}, \file{.sby}] The decode job. It compares the
  address decoder of each fabric that BAG generates with a reference that BAG writes from the IR.
\item[\file{formal/<project>_route_<unit>.sv}, \file{.sby}] One route job for each fabric
  and each bridge that BAG generates, named after its instance. It checks that unit on its own: routing,
  arbitration and bridge windows.
\item[\file{formal/bag_step.v}] Registers that the property files use to count solver
  steps.
\item[\file{formal/bag_faxi4.sv}, \file{formal/bag_fahb.sv}] BAG's own protocol checkers
  for AXI4 and AHB-Lite ports. BAG writes each one only when a job needs it.
\item[\file{formal/vendor/zipcpu/}] The protocol checkers from ZipCPU that the jobs use,
  for AXI4-Lite and APB ports, with their \file{LICENSE} and a \file{SOURCE.md} that
  names their source and commit.
\end{description}

Section~\ref{sec:formal} explains what each job checks and how to read its result. One
job, the quiet-in-reset job, is shown here:

\lstinputlisting[style=console]{examples/console/out-reset-job.txt}

\section{The file header}
\label{sec:header}

Every file that BAG generates from your spec starts with a header. The header is the
text of the file that Wrapper > \setting{File header} names, followed by BAG's banner. The banner names the version of BAG, the project, its revision, and the SHA-256
hash of the IR. It says \code{DO NOT EDIT}: change the spec and generate again. In SystemVerilog and C files the text is copied as it is, so write it as \code{//}
line comments: BAG does not check this, and other text stops the RTL from compiling. In
Markdown and HTML the whole header goes inside an HTML comment, and in Python and
\file{.sby} files BAG puts \texttt{\#} before each of its lines.

Three files have no header: the file list, the manifest and the copy of the IR in
\file{tests/}. The files that BAG copies (the in-house and vendored RTL, the ZipCPU
checkers, \file{bag_step.v}, \file{bag_faxi4.sv} and \file{bag_fahb.sv}) keep their own
headers.

\section{Same input, same output}
\label{sec:deterministic}

\cmd{bag generate} is \term{deterministic}: for the same IR, the same File header text
and the same version of BAG, it writes the same bytes on every run and on every machine.
No generated file holds a time stamp, a user name or a path of your machine. Every text
file has LF line ends, also on Windows. A workbook and the IR that \cmd{bag ir} writes
from it give the same output. So you can keep the output directory under version
control and review a change of the spec as a diff of the generated files.

The IR hash in the header changes when anything in the IR changes. The version in the
header changes when you update BAG, so a new version changes every generated file, even
when nothing else changed.

The folders that the EDA flows write (\file{obj/}, \file{sim/}, \file{sim_ready_shim/}
and the work folders of the formal jobs under \file{formal/}) are not part of this: the
tools write them, and their content can differ between runs.

\section{Which settings switch each part on}
\label{sec:switches}

The Project sheet decides which parts BAG writes. A setting is \code{Y} or \code{N}
(chapter~\ref{ch:workbook}).

\begin{center}
\begin{tabularx}{\linewidth}{>{\raggedright\arraybackslash}p{0.34\linewidth}X}
\toprule
Setting & What it switches \\
\midrule
\setting{Generate RTL} & \file{rtl/}, the file list, \file{doc/tieoffs.md}, \file{tests/}
  and \file{formal/}. With \code{N}, BAG writes only the C header and the other documents,
  and the manifest records no top, so lint, simulation and formal say no RTL was
  generated. \\
\setting{Generate C header} & \file{sw/<project>_memmap.h}. \\
\setting{Generate memory-map doc} & All of \file{doc/}: the memory map in Markdown and
  HTML, the port map, the reset sequence, and the tie-off report (with \setting{Generate RTL}
  = \code{Y}). \\
\setting{Generate cocotb unit tests} & All of \file{tests/}: the tests, the copy of the
  IR and the READY shim (with \setting{Generate RTL} = \code{Y}). \\
\setting{Generate formal (SymbiYosys)} & All of \file{formal/} (with \setting{Generate RTL}
  = \code{Y}). \\
\setting{Lint tools} & No file: the tools that \cmd{bag lint} runs, recorded in the
  manifest. \\
\setting{Simulator} & No file: the simulator of \cmd{bag sim}, recorded in the manifest.
  Only \code{verilator} is accepted. \\
\bottomrule
\end{tabularx}
\end{center}

With \setting{Generate RTL} = \code{Y}, the parameter package, the top wrapper, the
sub-wrappers, the copied RTL and the file list are always written; the manifest always
is. Wrapper > \setting{Port style} and Wrapper >
\setting{File header} change what the RTL looks like, but not which parts are written.

\section{A user-supplied fabric}
\label{sec:external}

A fabric whose template is \code{bag_external} is a \term{user-supplied fabric}: you
write its RTL, and BAG instantiates it in the place where a generated fabric would be.
Three things change in the output directory.

\begin{itemize}
\item BAG writes a \term{stub} for the fabric, \file{rtl/ext/<module>.sv}, and lists it in
  the file list. The module name is Fabrics > \col{External module} or, if that is
  blank, \code{<project>_<fabric>}. The stub has the exact ports that BAG connects. Its
  outputs are tied to zero, so the design lints and elaborates, but the stub accepts and
  answers nothing.
\item There are no generated tests: \file{tests/} is not written, and
  \cmd{bag generate} says so. The tests drive every fabric through the top, and BAG
  does not know how your fabric behaves.
\item There are no formal jobs: \file{formal/} is not written, for the same reason.
\end{itemize}

This holds for the whole design: one user-supplied fabric turns off the tests and the
formal jobs of every fabric in it. The C header and the documents are written as usual.

\lstinputlisting[style=console]{examples/console/out-external.txt}

The ports of the stub are, in order: \code{clk_i} and \code{rst_ni} (the fabric's clock
and its conditioned, active-low reset); the bus signals of each attached port, master
side first, named by Wrapper > \setting{Signal naming pattern}; then the ports that the
\sheet{User_Ports} sheet declares for the fabric. Directions are the fabric's own: on a master's
port the fabric is the slave end. Chapter~\ref{ch:workbook} describes how a shared
External module names its ports by position instead.

Replace the stub with your fabric's RTL. \cmd{bag generate} writes the stub again on
every run, so keep your RTL outside the output directory and copy it over the stub
after each run, or compile your own file in its place.
